\documentclass[../../../include/open-logic-section]{subfiles}

\begin{document}

\olfileid[id]{sth}{ord-arithmetic}{expo}
\olsection{Eksponensiasi Ordinal}

Sekarang kita beralih ke eksponensiasi ordinal. Sayangnya, tidak ada definisi
sintetis yang \emph{bagus} bagi eksponensiasi ordinal.

Memang, definisi sintetis yang eksplisit \emph{ada}. Berikut salah satunya.
Misalkan $\text{finfun}(\alpha,\beta)$ himpunan semua fungsi $f \colon
\beta \to \alpha$ sedemikian sehingga $\Setabs{\gamma \in
\beta}{f(\gamma) \neq 0}$ sama banyak dengan suatu bilangan asli.
Definisikan suatu pengurutan baik pada $\text{finfun}(\alpha,\beta)$ dengan
$f \sqsubset g$ jika dan hanya jika $f \neq g$ dan $f(\gamma_0) <
g(\gamma_0)$, dengan $\gamma_0 = \text{max}\Setabs{\gamma \in
\beta}{f(\gamma) \neq g(\gamma)}$. Kemudian kita dapat mendefinisikan
$\ordexpo{\alpha}{\beta}$ sebagai $\ordtype{\text{finfun}(\alpha, \beta),
\sqsubset}$. Potter menggunakan definisi eksplisit ini, lalu langsung
menjelaskan:
\begin{quote}
	Pemilihan pengurutan ini ditentukan semata-mata oleh keinginan kita untuk
	memperoleh definisi eksponensiasi ordinal yang memenuhi syarat rekursif
	yang sesuai\ldots, dan pengurutan ini jauh lebih sulit dibayangkan daripada
	jumlah terurut maupun hasil kali terurut.
	\citep[hlm.~199]{Potter2004}
\end{quote}
Tepat sekali. Kita menjelaskan penjumlahan sebagai ``suatu salinan $\alpha$
yang diikuti suatu salinan $\beta$'', dan perkalian sebagai ``suatu barisan
sepanjang $\beta$ yang terdiri atas salinan-salinan $\alpha$''. Namun, kita
tidak mempunyai uraian singkat yang menarik mengenai
$\text{finfun}(\alpha, \gamma)$. Karena itu, sebagai gantinya kita akan
memberikan definisi eksponensiasi ordinal hanya \emph{melalui} rekursi
transfinit, yakni:

\begin{defn}\ollabel{ordexporecursion}
\begin{align*}
	\ordexpo{\alpha}{0} &= 1\\
	\ordexpo{\alpha}{\beta\ordplus 1} &=\ordexpo{\alpha}{\beta} \ordtimes \alpha\\
	\ordexpo{\alpha}{\beta} &= \bigcup_{\delta < \beta}\ordexpo{\alpha}{\delta}& & \text{jika $\beta$ merupakan ordinal limit}
\end{align*}
\end{defn}

Jika kita bekerja \emph{sebagai} teoretikus himpunan, kita mungkin ingin
menyelidiki beberapa sifat eksponensiasi ordinal. Namun, tidak banyak lagi
yang perlu kita tambahkan selain fakta yang tidak mengejutkan bahwa
eksponensiasi ordinal tidak komutatif. Jadi $\ordexpo{2}{\omega} =
\bigcup_{\delta < \omega}\ordexpo{2}{\delta} = \omega$, sedangkan
$\ordexpo{\omega}{2} = \omega \ordtimes \omega$. Lagi pula, kita memang
tidak seharusnya \emph{mengharapkan} eksponensiasi bersifat komutatif, sebab
pada bilangan asli pun tidak demikian: $\ordexpo{2}{3} = 8 < 9 =
\ordexpo{3}{2}$.

\begin{prob}
Dengan menggunakan Induksi Transfinit, buktikan bahwa jika kita mendefinisikan
$\ordexpo{\alpha}{\beta} = \ordtype{\text{finfun}(\alpha, \beta),
\sqsubset}$, kita memperoleh persamaan-persamaan rekursi dalam
\olref[sth][ord-arithmetic][expo]{ordexporecursion}.
\end{prob}

\end{document}
